% !TeX root = ../sections/02_exercises.tex
\subsection{1-C-1.}
\begin{exercise}
    数列 $\{a_n\}$ についての以下の問いに答えよ。

    \begin{enumerate}
        \item
        \[
            \lim_{n\to\infty} a_n=\infty
        \]
        ならば
        \[
            \lim_{n\to\infty}
            \frac{a_1+a_2+\cdots+a_n}{n}
            =\infty
        \]
        であることを証明せよ。

        \item
        $\{a_n\}$ が有界であるが
        \[
            \left\{
                \frac{a_1+a_2+\cdots+a_n}{n}
            \right\}
        \]
        が発散するような例をひとつ与えよ。
    \end{enumerate}
\end{exercise}
\begin{background}
    この問題では，数列そのものの発散と，その算術平均の振る舞いを調べる。

    \begin{definition}[正の無限大への発散]
        数列 $\{a_n\}$ が正の無限大に発散するとは，
        任意の $M>0$ に対して，ある自然数 $N$ が存在して，
        $n\geq N$ ならば
        \[
            a_n>M
        \]
        が成り立つことをいう。
        このとき
        \[
            \lim_{n\to\infty}a_n=\infty
        \]
        と書く。
    \end{definition}

    \begin{definition}[有界な数列]
        数列 $\{a_n\}$ が有界であるとは，ある実数 $C>0$ が存在して，
        すべての $n\in\mathbb{N}$ に対して
        \[
            |a_n|\leq C
        \]
        が成り立つことをいう。
    \end{definition}

    \begin{remark}
    (1) では，$a_n$ が十分大きな番号以降で任意に大きくなることを用いて，
    平均
        \[
            \frac{a_1+\cdots+a_n}{n}
        \]
        も任意に大きくなることを示す。

        ただし，最初の有限個の項
        \[
            a_1,\ldots,a_{N-1}
        \]
        は負である可能性もあるため，その影響を後半の大きな項で打ち消す必要がある。
    \end{remark}

    \begin{lemma}[有限個の和の影響は平均で消える]
        実数 $c$ を固定するとき，
        \[
            \frac{c}{n}\to 0
        \]
        である。
    \end{lemma}

    \begin{remark}
    (2) では，有界な数列であっても，平均が必ず収束するとは限らないことを示す。
        そのためには，$1$ と $-1$ のような有界な値を，長いブロックごとに交互に並べる例を考えるとよい。
    \end{remark}
\end{background}
\begin{strategy}
(1) では，
    \[
        \frac{a_1+a_2+\cdots+a_n}{n}
    \]
    が正の無限大に発散することを示す。

    任意の $M>0$ を取る。
    仮定より
    \[
        a_n\to\infty
    \]
    であるから，十分大きな $k$ について
    \[
        a_k>2M
    \]
    が成り立つ。

    そこで，ある自然数 $N$ が存在して，
    $k\geq N$ ならば
    \[
        a_k>2M
    \]
    となる。

    このとき，和を
    \[
        a_1+\cdots+a_n
        =
        (a_1+\cdots+a_{N-1})
        +
        (a_N+\cdots+a_n)
    \]
    と分ける。

    後半部分ではすべての項が $2M$ より大きいので，
    \[
        a_N+\cdots+a_n>2M(n-N+1)
    \]
    と評価できる。

    前半部分
    \[
        a_1+\cdots+a_{N-1}
    \]
    は有限個の和なので，$n$ で割ると影響は小さくなる。
    したがって，$n$ をさらに十分大きくすれば
    \[
        \frac{a_1+\cdots+a_n}{n}>M
    \]
    が得られる。

    よって，
    \[
        \frac{a_1+\cdots+a_n}{n}\to\infty
    \]
    である。

    (2) では，平均が収束しない有界列の例を作る。
    例えば，$1$ と $-1$ を長さが倍々に増えるブロックで交互に並べる。

    具体的には，
    \[
        a_n=(-1)^m
        \qquad
        (2^m\leq n<2^{m+1})
    \]
    と定める。

    このとき，すべての $n$ に対して
    \[
        |a_n|=1
    \]
    であるから，$\{a_n\}$ は有界である。

    一方，ブロックの終わり
    \[
        n=2^{m+1}-1
    \]
    で平均を調べると，部分和は
    \[
        1-2+4-8+\cdots+(-1)^m2^m
    \]
    となる。

    したがって，平均は偶数番目のブロック終端ではおよそ $1/3$ に近づき，
    奇数番目のブロック終端ではおよそ $-1/3$ に近づく。

    つまり，平均列は異なる極限をもつ部分列を含むので収束しない。
    したがって，求める例が得られる。
\end{strategy}